% !PRINTED 284
\weberpage{284}

Then the entire group $L_7$ is represented thus:

\weq{9}{\tau^\varrho \chi^\lambda \omega^\mu \Theta^\nu, \quad \begin{aligned} \varrho &= 0, 1, 2, 3, 4, 5. 6\\ \lambda &= 0, 1, 2; \ \mu = 0, 1\\ \nu &= 0, 1, 2, 3. \end{aligned}} % [sic] full stop printed between 5 and 6

That all these elements (9) are different from one another follows simply from the fact that no power of $\tau$ whose exponent is not divisible by 7 can be contained in the octahedral group, because the octahedral group has no element of degree $7$.

The group (5) of degree $21$ contained in the group $L_7$ is then represented in the form $\tau^\varrho \chi^\lambda$.

In order to characterise the composition of the elements (9), it evidently suffices if for every $\varrho$ the elements

\weq{10}{\Theta\tau^\varrho, \ \omega\tau^\varrho, \ \chi\tau^\varrho}

are represented in the form (9), because thereby, together with the compositions in the octahedral group, the symbolic product of any two elements of the form (9) can again be represented in the form (9).

Now since

\weq{11}{\tau^\varrho = \begin{pmatrix} 1, & \varrho \\ 0, & 1 \end{pmatrix}.}

one first obtains very simply from (6)

\weq{12}{\chi\tau^\varrho = \tau^{4\varrho}\chi,}

and from this

\weq{13}{\chi^2\tau^\varrho = \tau^{2\varrho}\chi^2,}

and one sees easily that these six relations are a consequence of the single one:

\weq{14}{\chi\tau = \tau^4\chi.}

The remaining compositions (10), however, are obtained only by actual computation in the individual cases, in which the table (7) renders good service. One applies it in such a way that for each value $\varrho$ one forms the elements

\[
\tau^{\varrho'} \omega \tau^\varrho, \quad \tau^{\varrho'} \Theta \tau^\varrho
\]

and determines $\varrho'$ so that these elements are found in the table, which is always possible in only one way.
